\documentclass[11pt,a4paper]{article}
\usepackage[margin=25mm]{geometry}
\usepackage[T1]{fontenc}
\usepackage{lmodern,microtype,amsmath,amssymb,amsthm,mathtools,booktabs,array}
\usepackage{xcolor,mathrsfs}
\usepackage{needspace}
\usepackage[colorlinks=true,allcolors=blue!45!black]{hyperref}
\hypersetup{pdftitle={A seven-eighths theorem for Hecke characters over arbitrary number fields},pdfsubject={Research draft}}
\newtheorem{theorem}{Theorem}[section]
\newtheorem{proposition}[theorem]{Proposition}
\newtheorem{lemma}[theorem]{Lemma}
\newtheorem{corollary}[theorem]{Corollary}
\newtheorem{hypothesis}[theorem]{Analytic input}
\theoremstyle{remark}\newtheorem{remark}[theorem]{Remark}
\newcommand{\N}{\mathrm N}
\newcommand{\Rea}{\operatorname{Re}}
\newcommand{\eps}{\varepsilon}
\newcommand{\one}{\mathbf 1}
\newcommand{\dd}{\,d}
\newcommand{\bpsi}{\beta_\psi}
\newcommand{\Packet}{\mathscr F_\omega}
\newcommand{\OO}{\mathcal O}
\newcommand{\Xcal}{\mathcal X}
\newcommand{\Mcal}{\mathcal M}
\newcommand{\Qcal}{\mathcal Q}
\numberwithin{equation}{section}

\title{A seven-eighths theorem for Hecke characters\\
over arbitrary number fields\\[3pt]
\large Original source results, field transfer, and one closure for each fixed infinity type}
\author{AI-assisted research draft\thanks{AI assistance: OpenAI ChatGPT was used for mathematical exploration and drafting; OpenAI Codex was used for editing and compilation.}}
\date{8 October 2026}
\begin{document}
\maketitle
\begin{abstract}
Relative to the original theta constructions and canonical analytic
proofs in OpenAI's quasi-Riemann and primitive-roots papers, every
unitary Hecke character over every number field has a zero-free
half-plane $\Rea s>7/8$. We prove the necessary transfer to ideal
norms, unit quotients, and a fixed arbitrary infinity type. In
particular, the general-field marked inverse and paired fourth
moments are consequences of the cited constructions and the transfer
arguments supplied here, rather than additional hypotheses.
The proof passes from $F$ to $K=F(\zeta_{12})$ and works with
finite-order twists of one fixed character and its conjugate.
A single-row sixth-power amplification gives $11/12$ for this
family; the original dynamic fourth-moment capacity then closes
directly at $7/8$. The real losses and slot system are fixed before
the target conductor, while height orders and constants may depend
on it. The long original canonical proofs are retained as source
inputs and are not independently re-proved in full. No earlier
extension manuscript is used as a theorem source.
\end{abstract}

\section{Main theorem}

The two primary sources are OpenAI's original quasi-Riemann paper
\cite{Q}, denoted Q, and its primitive-roots paper \cite{P}, denoted P.
We distinguish the original analytic work that is cited from the
field and infinity-type transfer proved in this paper.

\begin{hypothesis}[Original source results]\label{hyp-source}
Retain the theta construction and local identities of P,
Sections 3 and 5--7 and Appendix A, and the physical estimates,
zero detection, and canonical moment arguments of Q,
Sections 4--18, with the finite-seminorm, strict-width, common-mask,
and polynomial-height assertions established there.
The full canonical constructions and reductions are used, not just
the published zero-free theorems. Their extension to a general field
and a fixed unitary infinity type is proved below.
\end{hypothesis}

The usual Hecke functional equation, class field theory, and the
number-field quadratic and higher-order large sieves enter through
the classical results used in Q and P. No earlier extension
manuscript is an additional input.

\begin{theorem}[Arbitrary number fields]\label{thm-main}
Under Source Input~\ref{hyp-source}, let $F$ be any number field.
Every continuous unitary Hecke character
$\psi:F^\times\backslash\mathbb A_F^\times\longrightarrow S^1$
satisfies
\begin{equation}\label{eq-main}
 \boxed{L_F(s,\psi)\ne0\qquad(\Rea s>7/8).}
\end{equation}
Finite and infinite order are both allowed, with arbitrary infinity
type. A pole for a unitary norm character is permitted. Additional
deleted Euler factors do not change the assertion. The boundary is
independent of the field, conductor, infinity type, and height;
implied constants and lower thresholds need not be uniform in these
data.

For a prescribed $\psi$, the proof uses only the fixed-type family
\eqref{eq-packet} over $K=F(\zeta_{12})$, with
$\omega=\psi\circ N_{K/F}$. It requires no estimate uniform over
all infinity types and no additional general-field moment hypothesis.
\end{theorem}

For a character in the normalization
$\Psi(\mathfrak a)=\psi(\mathfrak a)(\N\mathfrak a)^u$,
$u\in\mathbb R$, one has $L_F(s,\Psi)=L_F(s-u,\psi)$.
Thus the corresponding boundary is $u+7/8$. In particular,
an algebraic character with ideal-value absolute size
$(\N\mathfrak a)^{w/2}$ has boundary $w/2+7/8$ in its
unnormalized variable.

There is one analytic closure:
\[
 \text{marked inverse moment}\ \Longrightarrow\ 11/12
 \ \Longrightarrow\ \text{dynamic fourth moment and closure}
 \ \Longrightarrow\ 7/8.
\]
The transfer is established in Section~\ref{sec-inputs} and
Appendix~\ref{app-transfer}; the descent from $K$ to $F$ is
elementary after standard Hecke continuation. The source-relative
qualification concerns the original long analytic constructions,
not a separate assumption of their desired general-field conclusion.

\section{Field reduction and the fixed-type family}\label{sec-reduction}

\subsection{Abelian descent}

The extension $K=F(\zeta_{12})$ is abelian over $F$, has degree at
most four, and is totally imaginary. It need not be cyclotomic over
$\mathbb Q$. The distinction matters: the cyclotomic scope of P,
Theorem 1.2, does not by itself cover this $K$.

We use ideal-value normalization for norm twists:
$\N_K^{it}(\mathfrak a)=(\N\mathfrak a)^{it}$.
Such a twist replaces $L_K(s,\eta)$ by $L_K(s-it,\eta)$.

\begin{lemma}[Descent, including possible poles]\label{lem-descent}
Let $K/F$ be a finite abelian extension and let $\psi$ be a unitary
Hecke character of $F$. If $L_K(s,\psi\circ N_{K/F})$ has no zeros
in $\Rea s>\sigma_0$, where $\sigma_0>0$, then
$L_F(s,\psi)$ has no zeros there.
\end{lemma}

\begin{proof}
Let $G=\operatorname{Gal}(K/F)$, and regard $\widehat G$ as a group
of finite-order Hecke characters of $F$ by class field theory.
Choose a finite set $S$ containing all ramified primes and all
conductors involved. With the corresponding primes omitted upstairs,
\begin{equation}\label{eq-descent}
 L_K^S(s,\psi\circ N_{K/F})
       =\prod_{\chi\in\widehat G}L_F^S(s,\psi\chi).
\end{equation}
One can verify this directly on the absolute half-plane. If an
unramified prime $\mathfrak p$ has Frobenius order $f$, it has
$|G|/f$ primes above it, each of norm $(\N\mathfrak p)^f$.
The values $\chi(\operatorname{Frob}_{\mathfrak p})$ run through
the $f$-th roots of unity, each $|G|/f$ times. Hence the two Euler
factors agree. Meromorphic continuation gives \eqref{eq-descent}
everywhere. The omitted Euler polynomials are nonzero in
$\Rea s>0$.

If all factors on the right are holomorphic at $s_0$, a zero of any
one of them gives a zero on the left. The only remaining possibility
is a pole. By the standard Hecke pole classification \cite{Tate},
at most one factor has a pole at $s_0$, and then
$\psi\chi_0=\N_F^{it}$ and $s_0=1+it$. Its pullback is
$\N_K^{it}$, so the left side and the $\chi_0$ factor both have
a simple pole. Taking meromorphic orders in \eqref{eq-descent},
\[
 -1=-1+\sum_{\chi\ne\chi_0}
              \operatorname{ord}_{s_0}L_F^S(s,\psi\chi).
\]
Every term in this sum is nonnegative, hence zero. No zero of
another factor can cancel the pole. This proves the assertion
on the whole half-plane, including its intersection with $\Rea s=1$.
\end{proof}

\subsection{One family for each fixed infinity type}

We now fix a totally imaginary number field $K$ containing
$\mu_{12}$, and one unitary Hecke character $\omega$ of $K$.
Let
\begin{equation}\label{eq-packet}
 \Packet=
 \{\theta\,\omega^e\N_K^{it}:
          e\in\{-1,0,1\},\
          \theta\text{ finite order},\ t\in\mathbb R\}.
\end{equation}
All finite-order characters are included by $e=0$. This family is
closed under finite-order twisting, conjugation, and vertical
translation of its $L$-functions. These are exactly the operations
on live Hecke denominators in the argument. Extracted coefficients
such as $\omega(\mathfrak b)^3$ do not require a live background
$\omega^3$; Section \ref{sec-inverse} makes this distinction explicit.

The real parts of zeros are unchanged by the parameter $t$, so
it may be omitted when taking suprema. After the starting bound is
proved, set
\begin{equation}\label{eq-common-B}
 B=B_{K,\omega}:=
 \sup\left(\{1/2\}\cup
   \{\Rea\rho:\eta\in\Packet,\ L_K(\rho,\eta)=0,\
                            1/2\le\Rea\rho\le1\}\right).
\end{equation}
The contradiction will use this one supremum for both the target
and its companion denominators. It will never substitute a
target's individual zero supremum for the common family boundary.

\section{Transfer of the original constructions}\label{sec-inputs}

All sums below use ideal norms. We write $q_{\mathfrak a}=\N\mathfrak a$.
The symbols $u,n,p$ in the source-compatible formulas denote ideal
indices with the specified finite class and unit sectors, or their
allowed generators in the $S$-integer presentation of P, Section 3.
They do not mean all algebraic integers of bounded absolute norm.
The latter set can be infinite because of units.

Fix $S$ large enough to include the archimedean places, the primes
over $6$, the fixed finite data, and representatives that trivialize
the $S$-ideal class group. Use P's character-equivariant integration
and periodization modulo $U_S^6$. In particular, the continuous
unit phases of $\omega$ are retained. Only the finite power-class
and ray tables are subjected to finite Fourier inversion.
Zero extensions of residue symbols remain in every formula.

\begin{proposition}[Quantitative field and type transfer]\label{hyp-main}
Under Source Input~\ref{hyp-source}, fix $K\supset\mu_{12}$ and
a unitary Hecke character $\omega$. The following three conclusions
hold for the family \eqref{eq-packet}. We use $\mathsf T$, $\mathsf I$,
and $\mathsf P$ as names for these consequences.
\begin{enumerate}
\item[\(\mathsf T\).]
The normalized cubic theta and Poisson constructions have the
general-field form of P, Sections 3 and 5--7 and Appendix A,
and the fully marked physical probe has the norm estimates of
Q, Sections 5--16. This includes the actual unit-quotient profiles,
finite local scalars, whole-index masks, all principal and zero
orbits, and the full compensated Euler correction. More precisely,
it supplies the normalized functions \eqref{eq-signal},
the principal assertions \eqref{eq-principal-approx}--\eqref{eq-contraction},
the low estimate \eqref{eq-low}, and the row decomposition
\eqref{eq-E-general} with the principal, small-row, and absolute
large-row estimates specified in Section \ref{sec-close}.
The zero detector and witness construction of Q, Section 8,
also hold with these norm scales. This transfer includes neither an improved row count nor a new
zero-free boundary.

\item[\(\mathsf I\).]
The canonical marked inverse estimate \eqref{eq-MI} holds with
Q, Section 17's strict widths, allowed bounded coefficients,
common masks, moving-puncture capacity, and finite-seminorm
profile rules. Its row indexing admits sixth-power amplification
with the original zero extensions. In particular, the independent
scales in \eqref{eq-two-scale} are available; the amplification
proof below explains this consequence of the unmarked estimate
and ideal counting.

\item[\(\mathsf P\).]
The paired ordinary fourth moment \eqref{eq-fourth} holds for
$3/4\le\kappa\le5/6$, with capacity
$n_1+n_2+6\kappa z\le m+q$, the zero-slot endpoint,
and the exceptional-row and common-mask conventions in
Section \ref{sec-fourth}. Its prime-cancellation premise is only
that finite-order Hecke characters have the common boundary
$(1+\kappa)/2$. The corresponding reciprocal and prime bounds
are used in the quantitative form of Q, Section 4.
\end{enumerate}
Each statement includes its finite-seminorm and polynomial-height
assertions, with the real choices in the following paragraph.
\end{proposition}

\begin{proof}[Location of the proof]
Proposition~\ref{prop-theta-transfer} establishes the general-field
arithmetic calibration. Lemmas~\ref{lem-general-arch} and
\ref{lem-orbit-gram} give the exact fixed-type reflection, masked
Poisson formula, and full additive Gram estimate. The finite marked
reflection and the two canonical inductions are transferred in
Appendix~\ref{app-transfer}. Lemma~\ref{lem-directed} verifies the
permitted infinite-type backgrounds, and the proof of
Proposition~\ref{prop-fourth} checks ordinary reflection, prime
pairing, and centered exceptional cancellation. These arguments use
the full source reductions and supply every field-dependent
replacement; they do not infer the moments from a large sieve alone.
\end{proof}

\subsection{The order of uniformity}

On compact admissible ranges of real parameters, for every prescribed
positive real loss, the width decrements, slot mesh, and finite slot
system in Proposition~\ref{hyp-main} can be chosen depending on $K$, the fixed
infinity type of $\omega$, and those losses, before the later
finite-order target twist and its conductor are selected.
Adding finitely many excluded primes may change constants and
thresholds but not these real choices.
In the fourth moment the mesh is chosen before, and independently
of, the fixed number of slots. Its eventual seminorm orders and
constants may depend on that number.

Implied constants, finite seminorm and Sobolev orders, height degrees,
height ceilings, and lower thresholds may depend on the full fixed
target and $S$. The external tail orders are chosen after the
height exponent. The transfer proofs give no additional
degree-dependent or $S$-dependent power of the scale $Z$:
fixed-degree unit logarithms consume a prescribed small loss,
and higher-dimensional Poisson decay is obtained by increasing
fixed derivative orders. The bounded induction depth depends
on the strict real margins; finite sector multiplicities change
constants and thresholds. Appendix~\ref{app-transfer} records
the resulting order of choices.

Thus the argument does not require estimates uniform over all
infinity types, all fields, or all conductors. It requires a common
real saving within the later finite-order twists of one already
fixed infinity type. These are the quantifiers used in
Section \ref{sec-close}.

\subsection{General-field theta calibration}
\label{subsec-theta-transfer}

\begin{proposition}[Field transfer of the theta construction]
\label{prop-theta-transfer}
Fix $K\supset\mu_{12}$.  Retain the original cubic theta construction
and local normalizations of \cite[Proposition 5.1 and Appendix A]{P}.
Its arithmetic realization, Gaussian smoothing, annular reflection,
and local Euler algebra apply over $K$ to every fixed unitary target,
after the conjugation choice of Lemma~\ref{lem-general-arch}.
These are consequences of that source construction, rather than
additional general-field analytic assumptions.
\end{proposition}

\begin{proof}
Choose $S$ to contain the primes over $6$, the trace-character
ramification, the target's finite conductor, the source's fixed local
data, and primes generating the ideal class group.  Then
$R=\mathcal O_{K,S}$ is principal.  Adelic additive duality and the
product formula identify $R$ with its annihilator in $K_S$ and give
it covolume one in the self-dual normalization, exactly as in
\cite[Lemma 3.1]{P}; the argument does not use a cyclotomic extension.

Put $[K:\mathbb Q]=2r$ and $s_f=|S_f|$.  For $\ell=2,3$, local
Kummer theory and the $S$-unit theorem give
\[
 \dim_{\mathbb F_\ell}K_S^\times/K_S^{\times\ell}=2s_f+2r,
 \qquad
 \dim_{\mathbb F_\ell}R^\times/(R^\times)^\ell=s_f+r.
\]
The second space injects into the first: a unit locally an
$\ell$th power at $S$ defines a Kummer extension unramified outside
$S$ and split at $S$, hence an everywhere unramified abelian extension
classified by the trivial $S$-ideal class group.  Hilbert reciprocity
makes the image isotropic.  Also $-1=\zeta_{12}^6$ is a sixth power,
so the sextic Hilbert pairing is alternating.  The image is therefore
Lagrangian, and it has an isotropic complement $B$.  Combining the
$2$- and $3$-parts gives the allowed generators of
\cite[Lemma 3.2]{P}, unique modulo $\Gamma=(R^\times)^6$.
Compactness of $K_S^1/\Gamma$ supplies balanced representatives.

For the cubic Gauss correction set
$\Lambda_K=\Lambda_{\mathbb Q(\sqrt{-3})}\circ N_{K/\mathbb Q(\sqrt{-3})}$.
At a prime of residue degree $f$, Davenport--Hasse lifting gives
$J_{q^f}=(-1)^{f-1}J_q^f$; the prime Jacobi identity consequently gives
$\gamma_2(\mathfrak p)^3=-\Lambda_K(\mathfrak p)$.
Each complex coordinate of the relative norm contributes either $z$
or $\bar z$, so the angular weights of $\Lambda_K$ are $\pm1$.
The quadratic correction follows from the Weil-index product formula
on the same isotropic complement.  This proves the identities of
\cite[Lemma 3.4]{P} over $K$, without a Galois or CM assumption.
The Kazhdan--Patterson specialization used in \cite[Appendix A]{P}
is over a general number field containing the required roots of unity:
its $(2,3,2)$ scalar quotient and exceptional datum are local, while
rational compatibility follows from reciprocity and the product
formula.  An everywhere-local cube is a global cube by the Kummer
extension and Chebotarev argument given there.  Thus no cyclotomic
restriction enters this specialization of the retained construction.

For an infinite-order target, only its smooth local unit restrictions
are cut off at finite places.  They have open kernels even if their
uniformizer values have infinite order.  Lemma~5.3 of \cite{P}
therefore supplies a common cutoff inside these kernels and the
sixth powers.  The target itself is never expanded in a finite
Fourier table.  At infinity use Lemma~\ref{lem-general-arch}, keeping
all relative radial parameters.  Its exact orbit substitution puts
the whole target value on the completed ideal index, and its Mellin
calculation supplies the Gaussian and arbitrary annular tests.
In particular no extra norm phase is introduced in the physical
scale monomial.  The exterior residue fields have
$q_{\mathfrak p}\equiv1\pmod{12}$, so all remaining local digit,
Gauss, and Euler computations of \cite[Sections 6--7]{P} are unchanged.
\end{proof}


\subsection{The archimedean interface over a general field}
\label{subsec-general-arch}

Let $K\supset\mu_{12}$ be fixed, with $[K:\mathbb Q]=2r$.
Use the theta representation, rational splitting, and coefficient
normalizations constructed in the arithmetic part of
Proposition~\ref{prop-theta-transfer}.
The following local and norm-sheet calculation uses
\cite[Proposition 5.1 and Appendix A]{P}, retaining every fixed
angular and relative radial parameter.

Put $R=\mathcal O_{K,S}$ and $\Gamma=(R^\times)^6$, with $S$ enlarged
as in \cite[Section 3]{P}.  Let $\mathscr L$ be the actual unitary
$S$-period character, including the Gauss correction, so
$\mathscr L(\Gamma)=1$.  At a complex place write
\[
 \mathscr L_v(z)=\omega_{m_v}(z)|z|_v^{i\tau_v},\qquad
 \omega_m(z)=(z/|z|_{\rm ord})^m,
 \qquad m_v\in\mathbb Z,\quad\tau_v\in\mathbb R.
\]
The relative parameters $\tau_v$ are retained; a single global norm
twist need not remove them.  The Gauss correction has angular
exponents $b_v\in\{1,-1\}$ by \cite[Lemma 3.4]{P}.
In the ideal-character convention, the period weights are
$m_v=k_v-b_v$ for target angular weights $k_v$.
Replacing the target by its conjugate if necessary ensures
$m_v\ne0$ at at least one place: the two vectors
$(k_v-b_v)_v$ and $(-k_v-b_v)_v$ cannot both vanish.
Finite-order children do not change these weights.

\begin{lemma}[Fixed-type norm-sheet reflection]\label{lem-general-arch}
Assume the preceding nonzero angular condition.  The theta period has
an exact reflected transform for every smooth test supported in a
fixed compact subinterval of $(0,\infty)$.  Its archimedean transform,
including each fixed finite-$S$ branch, satisfies
\[
 |(x\partial_x)^h\Phi_{\kappa,V}(x)|
 \ll_{K,\mathscr L,A,h}\mathcal S_J(V)
       \min\{x^{1/4},x^{-A}\},
 \qquad J>4rA+h+1,
\]
where $\mathcal S_J$ is a finite smooth seminorm on that interval.
For the fixed Gaussian test the same Mellin identity applies;
its constants may depend on the fixed type.  A moving norm phase $V(x)x^{iT}$
costs at most $(1+|T|)^J$ in this estimate.
\end{lemma}

\begin{proof}
At a complex place let $\phi^\circ$ be spherical and choose the
additive coordinate as in \cite[Section 5.2]{P}.  The holomorphic
and antiholomorphic upper-root operators $E,\bar E$ can be normalized
with the same scalar so that they multiply the Whittaker function
by $z,\bar z$, respectively; this is (5.10) of \cite{P}.
Take $\phi_m=E^m\phi^\circ$ for $m\ge0$, and
$\phi_m=\bar E^{-m}\phi^\circ$ for $m<0$.
With $a_v=(|m_v|+1)/2$, the local Whittaker function is
\[
 W_{\phi_{m_v}}(a(z))
   =|z|_v^{a_v}\omega_{m_v}(z)
     K_{1/3}(c_v|z|_v^{1/2}),\qquad c_v>0,
\]
after a common spherical normalization.
The complexified compact Lie algebra annihilates $\phi^\circ$;
in particular $(F-\bar E)\phi^\circ=0$.
Since $F$ and $\bar E$ commute, $F^j\phi^\circ=\bar E^j\phi^\circ$.
Write $\mathcal K_v$ for the maximal compact subgroup. Then
$w\in\mathcal K_v$ and $\operatorname{Ad}(w)E=-F$, giving the exact identity
\[
                    \pi(w)\phi_m=(-1)^{|m|}\phi_{-m}.
\]
Thus each complex place is treated independently; no common angular
weight or single global derivative is used.

At the finite places choose the Whittaker cutoff $\one_U$ from
\cite[Lemma 5.3]{P}, with $U$ contained in the sixth powers and in
$\ker\mathscr L$.  Write
$X_u=\{y\in K_S^\times:|y|_S=u\}$ with the sheet measure of
\cite[(5.15)]{P}.  The unreflected profile is
\[
 V_0(u)=\int_{X_u}\one_U(y_f)
       \prod_{v\mid\infty}t_v^{a_v-i\tau_v}
                          K_{1/3}(c_v\sqrt{t_v})\,d_u^\times y,
 \qquad t_v=|y_v|_v.
\]
Using $\widehat f(s)=\int_0^\infty f(x)x^s\,dx/x$, absolute Fubini
and the Bessel integral give, initially for $\Re s>-1/3$,
\[
 \widehat V_0(s)=C_0 A_0^s
   \prod_{v\mid\infty}\prod_{\epsilon=\pm1}
      \Gamma(a_v+s-i\tau_v+\epsilon/6),
 \qquad C_0\ne0,\quad A_0=\prod_v4/c_v^2>0.
\]
For a reflected branch $\kappa$, the character in the profile is
$\mathscr L$, rather than $\mathscr L^{-1}$, and the vector is
$\pi(w)\phi$.  Consequently the exact Mellin multiplier is
\begin{equation}\label{eq-general-arch-ratio}
 \mathcal R_\kappa(s)
 :=\frac{\widehat H_\kappa(s)}{\widehat V_0(-s)}
 =C_\kappa A_0^{2s}B_\kappa(s)
  \prod_{v\mid\infty}\prod_{\epsilon=\pm1}
   \frac{\Gamma(a_v+s+i\tau_v+\epsilon/6)}
        {\Gamma(a_v-s-i\tau_v+\epsilon/6)} .
\end{equation}
Here $B_\kappa$ is a finite sum of finite-place Mellin factors.
For every $\delta<1/3$, (5.13) of \cite{P} makes it holomorphic
on $\Re s>-\delta$ and bounded on each closed vertical substrip.
The original finite factor is the nonzero constant $\operatorname{vol}(U)$;
no arbitrary finite Mellin factor is divided out.
Gamma poles have real part at most $-a_v+1/6\le-1/3$.
Taking $1/4<\delta<1/3$, the multiplier is therefore holomorphic
through the line $\Re s=-1/4$.  Stirling's formula cancels the
exponential factors separately at each place, since the imaginary
arguments in each quotient are opposite.  On a fixed vertical line
in this region,
\[
        |\mathcal R_\kappa(\sigma+it)|
             \ll_{K,\mathscr L,\sigma}(1+|t|)^{4r\sigma}.
\]

For completeness, this multiplier belongs to an exact period
identity, not merely a formal local transform.  One integrates on
$X_u/\Gamma$, which is compact by the $S$-unit theorem.
Rotating at a place with $m_v\ne0$ kills both constant terms,
because their inducing character is radial; this is the argument
of \cite[(6.8)]{P}.  The absolute bounds of
\cite[Lemma 5.4]{P} apply unchanged to the fixed unitary character.
Choose a measurable fundamental domain $\Omega_u\subset X_u$ for
$X_u/\Gamma$. Orbit unfolding is the identity
\[
 \sum_{\gamma\in\Gamma}\int_{\Omega_u}
 \one_B([y])\mathscr L(y)^{-1} f(\gamma Ay)\,d_u^\times y
 =\mathscr L(A)\int_{X_{q_Au}}\one_B([A]^{-1}[v])
                    \mathscr L(v)^{-1}f(v)\,d^\times v,
\]
as in \cite[(6.13)]{P}.  On the unreflected cutoff $[v]=1$
and $[A]\in B$, so the remaining class indicator is one.
Thus ideals are counted once, and no count of their infinitely many
unit associates is substituted for a norm-sheet integral.
The two expansions give the completed Mellin functional equation.
Its continuation is justified by the rapid decay at opposite ends
after constant-term removal; reciprocal gamma factors are entire.
The convergent expansions on opposite half-planes and
Phragmen--Lindelof then give polynomial growth on fixed strips.

Mellin inversion now yields the reflected kernel
\[
 \Phi_{\kappa,V}(x)=\frac1{2\pi i}\int_{(\sigma)}
       \widehat V(-s)\mathcal R_\kappa(s)x^{-s}\,ds.
\]
Shift to $\sigma=-1/4$ or $\sigma=A$ and integrate by parts $J$
times in the Mellin transform of $V$.  This proves the estimate.
Finally $\mathcal S_J(Vx^{iT})\ll(1+|T|)^J\mathcal S_J(V)$.
One keeps this fixed-type reflection and changes the test; one does
not rebuild Gaussian deconvolution kernels for each $T$.
\end{proof}

The field degree increases the fixed seminorm order $J$, but adds
no norm-length exponent to this lemma.  Once all internal seminorm
orders have been fixed, a height $T_1=Z^\tau$ with sufficiently small
fixed $\tau>0$ absorbs its polynomial cost; the external tail order
is chosen afterwards.  The Poisson and row estimates are transferred next, and the
canonical moment reductions are completed in
Appendix~\ref{app-transfer}. These use the same finite-seminorm
profile class as the reflection lemma.

\subsection{Poisson and additive correlation}
\begin{lemma}[Norm-form Poisson and the additive Gram estimate]\label{lem-orbit-gram}
Fix a number field $K\supset\mu_{12}$ and fixed arithmetic data as in
P, Section~3. Let $R=\mathcal O_{K,S}$, and choose a finite-index power
subgroup $\Gamma\subset R^{\times6}$ on which the relevant finite
characters are trivial. Put $T=K_S^1$. The masked Poisson operations
in Q may be performed on $\Gamma$-orbits with ideal norms. Their
scale factors and the additive Gram bound Q, (15.4), are unchanged.
More precisely, for the additive polynomial of Q, Section~6, with its
specified finite ray coefficients and a fixed compact additive row
shape, put $P_a=Y^2/Q\ge1$. Then
\[
 \sum_m |A_m(Y)|^2
 \ll_{K,S,\epsilon}(1+|\nu|)^J\frac QY
       \left(1+P_a^{1/6}+\frac{P_a^2}{Y}\right)Z^\epsilon.
\]
Here $\nu\in\mathbb R$ is the norm-power height in the additive test,
the row shape has norm scale $Q$, the scales lie in fixed
polynomial ranges, and $J$ is a finite fixed seminorm order.
No arbitrary row-dependent arithmetic coefficient is allowed.
The real loss $\epsilon$ is prescribed before the fixed target and
its finite conductor; constants, $J$, and thresholds may depend on
those fixed data.
\end{lemma}

\begin{proof}
The quotient $T/\Gamma$ is compact. Choose a smooth, compactly
supported unit partition on $T$, with the finite-place components
locally constant, whose translates by $\Gamma$ sum to one.
This represents orbit sums by lattice sums without counting infinitely
many associates. If $\ell_x$ is the positive norm section, define
\[
 W_\phi(x)=\int_T\phi(\ell_xv)\,d^\times v,
 \qquad \widetilde W_\phi(x)
       =\int_T\widehat\phi(\ell_xv)\,d^\times v.
\]
Every compact smooth annular norm test has such a lift, by multiplying
it by a fixed compact shape function of integral one. For a primitive
exterior residue character $\chi$ of conductor $\mathfrak m=mR$,
the nonprincipal masked formula is
\[
 \sum_{[k]\ne0}\chi(k)\one_{(k,\mathfrak R)=1}W_\phi(q_k/L)
 =\frac{L\gamma(\chi)}{\sqrt{\N\mathfrak m}}
  \sum_{\mathfrak d\mid\operatorname{rad}\mathfrak R}
       \frac{\mu(\mathfrak d)\chi(d)}{\N\mathfrak d}
  \sum_{[h]\ne0}\bar\chi(h)
       \widetilde W_\phi\!\left(
         \frac{Lq_h}{\N\mathfrak d\,\N\mathfrak m}\right).
\]
Indeed, expand the mask, substitute $k=dk'$, and apply additive
Poisson on the self-dual $R$-lattice while integrating on the compact
quotient $T/\Gamma$. The finite transform is the ordinary primitive
Gauss transform. Haar translation removes the norm-one part of $dm$,
so the transformed kernel is common and has no moving modulus shape.
Fixed local sectors are separated before this operation. For a general
periodic row coefficient one retains its full finite Fourier transform
instead of replacing it by a primitive character. The dual zero is
kept separately; for the primitive principal character it is
\[
 L\operatorname{vol}(T/\Gamma)\widehat\phi(0)
       \sum_{\mathfrak d\mid\operatorname{rad}\mathfrak R}
                         \frac{\mu(\mathfrak d)}{\N\mathfrak d}.
\]
Any included primal zero is the separate compact-quotient scalar
$\operatorname{vol}(T/\Gamma)\phi(0)$ with its actual mask.
It is never interpreted as $W_\phi(0)$.

\paragraph{The general periodic version, including both zero orbits.}
Let $f:R/\mathfrak c\to\mathbb C$ be a periodic function such that
$f(\gamma a)=f(a)$ for every $\gamma\in\Gamma$; every fixed local
sector and every original zero extension is retained in $f$ or in
the fixed local profile. For $\mathfrak d\mid\operatorname{rad}\mathfrak R$
define the unnormalized finite transform
\[
 \widehat f_d(h)=\sum_{a\bmod\mathfrak c}
           f(da)e(ha/c),\qquad\mathfrak c=cR.
\]
Use the additive-sign convention consistent with the chosen lattice
Fourier transform. Here $[k],[h]\in\Gamma\backslash(R\setminus\{0\})$,
$q_k=\N(kR)=|k|_S$, and $c,d$ are chosen exterior generators.
The exact orbit formula is
\[
 \sum_{[k]\ne0}f(k)\one_{(k,\mathfrak R)=1}W_\phi(q_k/L)
 =\frac L{q_c}\sum_{\mathfrak d\mid\operatorname{rad}\mathfrak R}
       \frac{\mu(\mathfrak d)}{q_d}
       \sum_{[h]\ne0}\widehat f_d(h)
           \widetilde W_\phi\!\left(\frac{Lq_h}{q_dq_c}\right)
       +\mathcal Z_f,
\]
where
\[
 \mathcal Z_f=
 \frac{L\operatorname{vol}(T/\Gamma)\widehat\phi(0)}{q_c}
       \sum_{\mathfrak d\mid\operatorname{rad}\mathfrak R}
              \frac{\mu(\mathfrak d)}{q_d}\widehat f_d(0)
 -\operatorname{vol}(T/\Gamma)\phi(0)f(0)
       \sum_{\mathfrak d\mid\operatorname{rad}\mathfrak R}\mu(\mathfrak d).
\]
The last term is absent for an annular lift. The mask is expanded
before the finite transform: for general $f$ one transforms $f(da)$,
not a fictitious multiplicative factor $f(d)f(a)$. Since every
$\gamma\in\Gamma$ is a unit modulo the exterior modulus, a residue
substitution shows that $\widehat f_d$ is also $\Gamma$-invariant.
Both nonzero sums are therefore genuine orbit sums. Specializing to
a primitive character gives
$\widehat f_d(h)=\chi(d)\sqrt{q_c}\gamma(\chi)\bar\chi(h)$
and recovers the preceding character formula. For the primitive
principal character $\mathfrak c=R$, the dual-zero formula specializes
to the displayed principal scalar.


On positive norm sheets, product Schwartz bounds give arbitrary
power decay at infinity and at most $(1+|\log x|)^b$ at zero,
where $b$ depends only on the fixed field and finite data.
Euler derivatives have the same logarithmic-degree bound. Thus the
additional logarithms consume a prescribed small real loss, with only
the threshold depending on the data. A positive Gaussian/compact-open
lift majorizes the entire row ball.

For the Gram bound follow the full-correlation proof of Q,
Proposition~15.2, retaining its collision zeros. Write
$s_i=Cn_i$, $(n_1,n_2)=1$. The first Poisson transform has precisely
the prefactor $Q/Y^3$: the lattice denominator $q_{s_1}q_{s_2}$
cancels the same factor in the complete Gauss correlation.
The zero frequency forces $n_1=n_2=1$ and contributes $O(Q/Y)$.
Insert the complete divisor identity for coprimality before extending
the coefficient, and write $n_i=dm_i$. Its bounded extension is used
with all original shared-prime zeros. After reciprocity the exterior
period is
\[
 \mathfrak r=\operatorname{lcm}(\operatorname{rad} C,
                 \operatorname{rad}(k)),\qquad
 q_{\mathfrak r}\ll_{\rm fixed}q_Cq_k,
\]
with the fixed local data included in the fixed profile or sectors.
For a nonexceptional $k$, namely one with a valuation not divisible
by six at a prime outside $C$ and the fixed support, the normalized
two-variable finite transform has size at most $q_C$ and its zero
coefficient vanishes. These are the same finite-field and CRT
identities as in Q; the residue-field cardinality simply becomes
$q_{\mathfrak p}$.

Keep one joint two-column profile, rather than separately chosen
row-dependent Fourier densities. On a shell
$q_Cq_k/P_a\asymp\mathcal R$, this profile has homogeneous seminorms
$O((1+|\nu|)^J\mathcal R^{-B})$. Split $q_C\asymp C_0$ and $q_d\asymp D_0$ into dyads.
With $N=Y/(C_0D_0)$, use the unit partition in both columns and balance an
allowed generator of $\mathfrak r$ by P, Lemma~3.2. Its normalized
shape stays in a fixed compact set. Two-column lattice Poisson gives
\[
 \left|\sum_{m\in R^2}a_{C,d,k}(m)W_{\mathcal R}(m/\ell_N)\right|
 \ll (1+|\nu|)^J\mathcal R^{-B}q_CN^2
          \min\{1,(q_{\mathfrak r}/N)^A\}.
\]
For $N\ge q_{\mathfrak r}$ the zero frequency is absent. A nonzero
dual vector has a large archimedean coordinate; taking the fixed
Schwartz order sufficiently large in terms of $A[K:\mathbb Q]$
gives the displayed decay. For $N<q_{\mathfrak r}$ use the compact
profile point count; a nonempty subunit annulus is bounded below by
a fixed constant. Thus field degree changes the derivative order,
not the exponent of $N$.

Ideal counting bounds the nonexceptional frequencies on this shell
by $O(\mathcal R P_a/C_0)$, up to a fixed-degree logarithm for
representatives in compact additive shapes. The same positive dyadic
sum as in Q is
\[
 \sum_{i,j\ge0}2^{-i-j}
       \min\{1,(2^{i+j}/\Lambda)^A\}
 \ll_A\frac{1+\log(2+\Lambda)}{\Lambda},
 \qquad\Lambda=Y/(\mathcal R P_a).
\]
After the shell sum this gives $(Q/Y)(P_a^2/Y)Z^\epsilon$.
Exceptional nonzero frequency ideals have the form
$\mathfrak a^6\mathfrak e$, with $\mathfrak e$ supported on $C$
and the fixed support. There are $q_C^\epsilon$ possibilities for
$\mathfrak e$ and $O((\mathcal R P_a/C_0)^{1/6})$ possibilities
for $\mathfrak a$. Representatives of one exterior ideal in the
compact additive shapes are $S$-units in a logarithmic box, hence
contribute only a fixed-degree logarithm. The corresponding term
is $(Q/Y)P_a^{1/6}Z^\epsilon$. Adding the diagonal gives the claim.
\end{proof}

\paragraph{Application to a fixed infinite type.}
In the physical probe the target character acts on the completed
theta column $cn^3$, not on the additive row variable $m$ or the
additive polynomial $A_m$. Consequently the preceding Gram argument
is literally unchanged for a genuine unitary target. Its continuous
infinity phase is handled by the fixed-type theta reflection lemma;
compensation contributes only the bounded column/slot coefficient
$\overline{\psi(p)}$. The completed-index masks, finite correlations,
local Euler algebra and subset scale coefficients remain the same.
The arbitrary-annular reflection has the same two-ended kernel
bound, with higher fixed seminorm and height orders. Therefore the
homogeneous common-profile separation and the positive row-energy
calculation of Q, Sections~14--15, retain their norm exponents.
At the original geometry the low estimate is $Z^{3/16+\epsilon}$.
The high Euler factors use only unit-modulus local target values;
their prime sums converge in the same regions over a fixed number
field. The principal residue constant is nonzero and may be calibrated
by one fixed scalar, as in P, Section~8. None of these steps imposes
a new zero-free hypothesis.


\section{The canonical family and the starting bound}\label{sec-inverse}

Fix a finite group $\Theta$ of finite-order ray characters containing
the fixed finite Gauss and reciprocity corrections. Enlarge the fixed set
$S$ to contain every fixed conductor and excluded prime. For a fixed target $\psi\in\Packet$, its genuine infinite-order
part is not made a member of $\Theta$. Global norm twists are placed
in the tests, and the finite-order cases use the original finite class.
For a finite-order target, enlarge $\Theta$ to contain $\psi$
and absorb $\psi$ into the finite presentation; in the paired
moment this case may always be written with $\omega_0=1$.
Writing $\chi_u(n)=\chi_n(u)$ for the physical sextic presentation,
the actual row family is
\begin{equation}\label{eq-family}
 \Xcal_u=
 \{\theta\chi_u,\theta\bar\chi_u:\theta\in\Theta\}
 \cup
 \{\psi\theta\bar\chi_u,\bar\psi\theta\chi_u:
                                     \theta\in\Theta\}.
\end{equation}
It is closed under conjugation. For each physical row it contains
finitely many characters, although some have infinite order.

\subsection{Preservation of the fixed infinity type}

For a presentation $\nu\chi_u^{\epsilon_\chi}$, the initialization
in Q, (17.68), conjugates the whole polynomial when
$\epsilon_\chi=+1$. Its canonical background is thus
\begin{equation}\label{eq-direction}
 \nu_{\rm can}=\nu^{-\epsilon_\chi}.
\end{equation}
The two conjugate genuine presentations in \eqref{eq-family}
both give a background in $\psi\Theta$; finite-order presentations
give one in $\Theta$.

\begin{lemma}[Directed canonical reduction]\label{lem-directed}
The canonical reductions of Q, Section 17, preserve the class
\[
 \nu_{\rm can}\in\Theta\cup\psi\Theta.
\]
No live background with a new power of the genuine character
is introduced. Together with the character-equivariant theta reflection and the
canonical transfer in Appendix~\ref{app-transfer}, this is the entire
infinite-type family needed by the marked inverse estimate.
\end{lemma}

\begin{proof}
We check the places where finite order could matter. In the first
Poisson step, Q, (17.22)--(17.30), write
$n_i=A_i C u_i$. On the original support,
\[
 \nu(n_1)\overline{\nu(n_2)}
 =\nu(A_1)\overline{\nu(A_2)}
       \nu(u_1)\overline{\nu(u_2)}.
\]
The common factor cancels because $|\nu(C)|=1$. Keep the last two
factors outside the finite Gauss-ray table. Expanding that table
then gives backgrounds $\nu\vartheta_i$, $\vartheta_i\in\Theta$;
no finite Fourier expansion of $\nu$ is being used. In the second
Poisson step, Q, (17.40)--(17.55), the common extracted factor again
contributes $|\nu\vartheta_i(v')|^2=1$. A child background is
$\nu\vartheta_i\vartheta$, still in the same class.

Cube extraction in Q, (14.2) and (17.10), gives $\nu(b)^3$
on the cube variable. This is an allowed bounded coefficient on
that variable, not a replacement of the live residual character
by $\nu^3$. Divisor extraction and the priority identities for
marks keep their original masks and product coefficients.

The remaining use of the background is the completed theta
terminal in Q, (14.14). Over $K$ this is supplied by
Lemma~\ref{lem-general-arch} and the finite marked reflection in
Appendix~\ref{app-transfer}. The finite table is kept separate from the fixed
archimedean phase. Its decay and finite-seminorm
bounds are exactly those used in the transferred canonical proof.
Thus the finite tables, masks, and live background are all
preserved, including at the reflected terminal.
\end{proof}

The resulting inverse estimate $\mathsf I$ has the following precise form. With the
source-compatible prime slots $Q_{u,i}$,
total slot length $z$, and a fixed polynomial-length inverse sum
\[
 \Mcal_u(D;W)=D^{-1/2}\sum_n
       \mu(n)\nu(n)\chi_n(u)^{\epsilon_\chi}W(q_n/D),
\]
take only presentations from \eqref{eq-family}, equivalently
$\nu^{-\epsilon_\chi}\in\Theta\cup\psi\Theta$.
The resulting estimate is
\begin{equation}\label{eq-MI}
 \sum_{0<q_u\ll Z^m}
 \left|\Mcal_u(Z^r;W)\prod_iQ_{u,i}\right|^2
 \ll Z^{m+\eps},
 \qquad r+2z<m,\quad 2r+8z<3m.
\end{equation}
Each strict inequality has a fixed positive margin. Slot coefficients
are bounded and fixed independently of the row; supports are
disjoint before masks are imposed. All products retain the same
coprimality mask. Only the rowwise profile choices permitted by Q,
Lemma 4.5, are allowed. At initialization the masks come only from
the fixed set $S$ and the displayed presentation zeros. Moving
punctures generated inside the canonical reduction must satisfy
Q, (17.2), with their capacity counted there. An additional freely
chosen moving mask is not allowed, even if held fixed during one
row sum.

\subsection{A single-row amplification}

\begin{proposition}\label{prop-seed}
Under Source Input~\ref{hyp-source}, the transferred inverse estimate implies,
for every $\eta\in\Packet$,
\begin{equation}\label{eq-seed}
 L_K(s,\eta)\ne0\qquad(\Rea s>11/12).
\end{equation}
This implication uses no prior zero-free bound for genuine
infinite-order characters.
\end{proposition}

\begin{proof}
Remove a global norm twist and use the permitted orientation
of the fixed infinite type; finite-order cases are included. Use $\nu=\eta$ and $\epsilon_\chi=-1$ in the inverse
polynomials. The first two-scale statement of Q, Lemma 17.6, applies
to independent $U,D\ge1$ and gives
\begin{equation}\label{eq-two-scale}
 \sum_{u\sim U}^{\text{sixth-power-free}}
 |\Mcal_u(D;W)|^2
 \ll \max\{U,U^{1/6}D^{5(1+c)/6}\}(UD)^\eps
 \qquad(c>0).
\end{equation}
Take a fixed bounded row shell containing $u=1$. This is the
independent-scale assertion, not the later specialization in which
$D=U^r$ and $r$ is bounded.

For completeness, the amplification in this special case is short.
Put $H=D^{1+c}$ and $P=H^{1/6}$. The no-slot case of
\eqref{eq-MI}, with base $H$, bounds the row energy up to $H$
by $H(HD)^\eps$, also after the scale supremum in Q, Lemma 4.5.
For each allowed ideal $a$ with $q_a\le P$, the exact identity is
\begin{equation}\label{eq-amplification-mask}
 \Mcal_1(D;W)=
 \sum_{d\mid\operatorname{rad}(a)}
       \mu(d)\eta(d)q_d^{-1/2}\Mcal_{a^6}(D/q_d;W).
\end{equation}
It follows from $\chi_n(a^6)=\one_{(n,a)=1}$.
There are $\gg_{K,S}P$ integral ideals $a$ in an allowed
fixed class with $q_a\le P$. Their sixth powers are distinct
ideal rows of norm at most $H$, and the divisor coefficient
mass is a small power. Weighted Cauchy and averaging give
\[
 |\Mcal_1(D;W)|^2\ll H/P\cdot D^\eps
             =D^{5(1+c)/6+\eps}.
\]
Letting $c$ and the preliminary loss be sufficiently small proves
\begin{equation}\label{eq-mobius}
 \sum_{(n,S)=1}\mu(n)\eta(n)W(q_n/D)
       \ll_{\eta,W,\eps}D^{11/12+\eps}.
\end{equation}

Choose a smooth dyadic partition $\sum_{j\ge0}W(q/2^j)=1$
for $q\ge1$. On a compact subset of $\Rea s>11/12$, the tests
$W_s(y)=W(y)y^{-s}$ have bounded relevant seminorms. Thus
\[
 G_\eta(s)=\sum_{j\ge0}2^{-js}
     \sum_{(n,S)=1}\mu(n)\eta(n)W_s(q_n/2^j)
\]
converges locally uniformly there by \eqref{eq-mobius}.
For $\Rea s>1$ it equals $1/L_K^S(s,\eta)$. The meromorphic
identity $L_K^S(s,\eta)G_\eta(s)=1$ proves the claim. Deleted
Euler factors have no zeros in $\Rea s>0$.

For two varying scales the same proof averages over $(u,a)$:
valuation modulo six recovers the sixth-power-free ideal $u$
from $ua^6$, and the quotient recovers $a$. The finitely many
class and unit sectors cost only a fixed factor. This proves
\eqref{eq-two-scale} from the same transferred row-energy estimate.

The starting argument uses only
$\mathsf T$ and $\mathsf I$, corresponding to Q, Section 17
and its theta terminals. It uses neither the fourth moment of Section 18
nor prime cancellation from a genuine-character zero-free region.
\end{proof}

\section{The paired fourth moment}\label{sec-fourth}

Proposition~\ref{prop-seed} gives $B\le11/12$ in
\eqref{eq-common-B}. In the contradiction $B>7/8$, put
\begin{equation}\label{eq-kappa}
 \Delta=B-\frac78,\qquad
 \kappa=2B-1=\frac34+2\Delta,\qquad
 0<\Delta\le\frac1{24}.
\end{equation}
Every finite-order Hecke character lies in $\Packet$.
Consequently the prime premise of the transferred fourth moment is already supplied
by the definition of $B$. A finite-order $7/8$ theorem is not
being assumed, and $\kappa$ is not replaced by $3/4$ in the
actual fourth-moment capacity.

Let $\xi_u$ be the finite moving presentation in Q, Section 18,
including its fixed finite-ray corrections and all zero extensions.
Let $\omega_0$ be either $1$, $\psi$, or $\bar\psi$, and put
$\lambda_u=\omega_0\xi_u$. For a positive-length prime slot,
take the paired coefficient $\omega_0^{-1}\vartheta_i$,
$\vartheta_i\in\Theta$. Its actual prime polynomial is
\begin{equation}\label{eq-paired}
 Q_{u,i}=P_i^{-1/2}\sum_p
       \xi_u(p)\vartheta_i(p)W_i(q_p/P_i).
\end{equation}
Consequently the character that needs cancellation at primes
is finite order.

\begin{proposition}[Paired fourth moment]\label{prop-fourth}
Under Source Input~\ref{hyp-source}, retain the row, common-mask, and fixed
moving-support hypotheses
of Q, Lemma 18.1, and write its effective width as $M=m+q$.
For a sufficiently fine fixed slot mesh,
\begin{equation}\label{eq-fourth}
 \sum_{u\in\mathcal R_z}
  |S_{\lambda_u}(n_1;W_1)S_{\lambda_u}(n_2;W_2)
                         \prod_iQ_{u,i}|^2
 \ll Z^{M+\eps}.
\end{equation}
Here $S$ denotes the source's centrally normalized plain sum.
For positive total slot length $z$, it suffices that
\[
 n_1+n_2+6\kappa z\le M,
\]
and $\mathcal R_z$ excludes rows whose finite inducing character
$\xi_u^*$ belongs to $\Theta$. For $z=0$, the bounded
nonnegative lengths $n_1,n_2$ are unrestricted, and only a
principal actual inducing character $\lambda_u^*$ is excluded.
The constants have the source's finite-seminorm and fixed
polynomial-height dependence.
\end{proposition}

\begin{proof}[Compatibility checks]
We record the changes in the source's fourth-moment reduction.
The canonical recursion and its field-dependent counting steps
are transferred in Appendix~\ref{app-transfer}. The following
checks supply its fixed-type reflection, paired prime premise,
and centered exceptional terminal.

First, ordinary Hecke reflection is available for $\lambda_u$.
Writing its fixed complex local parameters as $(k_v,\tau_v)$,
the gamma quotient is, up to fixed exponential scales,
\[
 \prod_{v\mid\infty}
 \frac{\Gamma(s+|k_v|/2+i\tau_v)}
      {\Gamma(1-s+|k_v|/2-i\tau_v)}.
\]
Its poles have nonpositive real part and it has polynomial growth
on fixed vertical strips away from its poles. The finite conductor
factor is $Q^{1/2-s}$, with $Q=|d_K|\N\mathfrak f_{\lambda_u}$,
so the dual norm length is still $Q/X$ \cite{Tate}.
Coincident poles can produce powers of $\log x$ at the small end
of the reflected profile. The central annular coefficients are
therefore bounded by $2^{-j/2}(1+j)^A$, still summable.
The common-mask identities Q, (18.4)--(18.8), must be applied
before a primitive inducing character is substituted.

Second, both finite Fourier transformations in Q,
Sections 18.4--18.6, act on the finite row character. Keep the
factors $\omega_0(n_1)\overline{\omega_0(n_2)}$ on the two
columns outside each finite table. Extraction multiplies them
by unit-modulus scalars; it leaves a single $\omega_0$ on
each surviving plain character. Its matching inverse remains
on the slot coefficient. Conjugating a reflected expression
is handled in two ways. After reflecting a single plain factor,
conjugate that factor and its test to recover the original
character inside its absolute square. A change of the whole
row orientation conjugates both plain factors and every slot,
including their coefficients and tests. Thus every live prime
polynomial continues to have a finite actual character as in
\eqref{eq-paired}. Q's prime estimate uses the common boundary $B$ and hence
$\kappa=2B-1$. The transferred norm-form capacity is
$n_1+n_2+6\kappa z\le M$. The dependence on $\Delta$ must
be retained in the later row comparison.

Third, the exceptional finite inducing family is handled on
unmasked functions before inserting a moving radical.
For each fixed $\vartheta\in\Theta$, Hecke continuation and a
Mellin shift give, uniformly for $X>0$,
\[
 \sum_n(\omega_0\vartheta)(n)W(q_n/X)(q_n/X)^{it}
 =c_{\omega_0\vartheta}\,X I_W(t)
       +O\bigl(p_J(W)(1+|t|)^A\bigr).
\]
After global norm twists have been removed, the coefficient
$c_{\omega_0\vartheta}$ is zero unless the character is principal;
$I_W(t)=\int_0^\infty W(y)y^{it}\,dy$.
For $X\ge1$ the error follows by shifting the unmasked $L$-function
to $\Rea s=-1$ and is in fact $O(X^{-1})$ with the same type
of seminorm bound. For $X<1$ bounded ideal counting suffices.

Insert the common radical $R$ only afterwards by
inclusion--exclusion. The main coefficient is multiplied by
$\sum_{d\mid R}\mu(d)(\omega_0\vartheta)(d)/q_d$,
independently of the scale; the error grows by at most $Z^\eps$
when $q_R$ is polynomially bounded. Keep the common kernel and
mask through the two equal-product rectangles before separating
the centered expression. Their principal main terms cancel
exactly as in Q, Section 18.7. For a genuine infinite-type
$\omega_0$ the residue is zero, giving a stronger bound.
Thus this exceptional calculation adds no degree-dependent
power of the moving norm. It preserves the exceptional-row
rule stated above.
More quantitatively, if the four scales in the two rectangles
are at least $Z^\varrho$, $\varrho\ge0$, and have common
product $T$, the unnormalized error after cancellation is
$O(Z^\eps T Z^{-\varrho}(1+|t|)^A)$.
Central normalization divides by $T^{1/2}$, giving
$O(Z^\eps T^{1/2}Z^{-\varrho}(1+|t|)^A)$, the saving
used in Q, (18.47).
\end{proof}

\section{Compensation and the principal factor}\label{sec-probe}

Let $I_{\psi;D}$ denote the physical probe in Q with target
$\psi$ and the mark $\one_{D\mid cn^3}$ on its whole
completed index. Proposition~\ref{prop-theta-transfer} supplies its character-equivariant
norm-form realization over $K$. At a good prime $p$,
$Q=q_p$, the compensation is
\begin{equation}\label{eq-comp-one}
 \overline{\psi(p)}I_{\psi;p}(X,Y,ZQ)
       -Q^{-3/2}I_\psi(X/Q,Y/Q,Z).
\end{equation}
For disjoint prime windows $\mathcal P_i$ and
$p_J=\prod_{i\in J}p_i$, its full form is
\begin{align}\label{eq-comp-full}
 I_{\psi,\mathrm{mod}}(Z)
 =\sum_{(p_i)}\prod_iW_i(q_{p_i}/P_i)
 \sum_{J\subseteq\{1,\ldots,J_0\}}&
       (-1)^{|J|}q_{p_J}^{-3/2}\overline{\psi(p_{J^c})}
       \nonumber\\[-2pt]
 &\hspace{-30mm}\cdot I_{\psi;p_{J^c}}
       (X/q_{p_J},Y/q_{p_J},Zq_{p_{J^c}}).
\end{align}
The target still acts on the whole completed index. All extra
target factors have absolute value one, so the subset rescaling
coefficients have exactly the sizes used by the low estimate
in $\mathsf T$.

The physical Mellin monomial is
$X^{1/2-z}Z^{s+z-1}Y^{w-1}$. Marking at $p$ and replacing
$Z$ by $ZQ$ contributes $Q^{s+z-1}P_p^*$; the term with
coefficient $Q^{-3/2}$ and scales $(X/Q,Y/Q)$ contributes
$Q^{z-w-1}P_p$. After extracting $Q^{z-1}$, their difference is
$\bar\psi(p)Q^sP_p^*-Q^{-w}P_p$. This verifies the scale
and phase in \eqref{eq-comp-one}.

For the remaining local calculation put
\[
 V=Q^{-6z},\quad W=\chi_u(p)Q^{-w},\quad
 D=\psi(p)\bar\chi_u(p)Q^{-s}.
\]
If $P_p$ is the local sum and $P_p^*$ its marked part, then
\[
 P_p-P_p^*=(1-V)^{-1}+W(1-W)^{-1},\qquad
 B_p=\overline{\psi(p)}Q^sP_p^*/P_p-Q^{-w}.
\]
Since $\overline{\psi(p)}Q^sD=\bar\chi_u(p)$,
the main prime factor is
\begin{equation}\label{eq-main-prime}
 \Qcal_i(u;z)=-\sum_{p\in\mathcal P_i}
       W_i(q_p/P_i)q_p^{z-1}\bar\chi_u(p).
\end{equation}
This is exactly the finite actual prime character needed in
Proposition \ref{prop-fourth}. No condition $\psi(p)=1$
is imposed on the prime windows.

The moment pairing can be made explicit. Orient a genuine
detector witness as $\lambda_u=\psi\theta\bar\chi_u$.
For the inverse moment, use direction $-1$ and relative slot
coefficient $(\psi\theta)^{-1}$, which is bounded and independent
of the row. For the plain fourth moment, conjugate both plain
witnesses and all selected prime slots together. The resulting
plain character is $\bar\psi\bar\theta\chi_u$, and its
relative slot coefficient is $\psi\theta$. Thus in Proposition
\ref{prop-fourth} one has $\omega_0=\bar\psi$,
$\xi_u=\bar\theta\chi_u$, and the required coefficient
$\omega_0^{-1}\theta$. The actual slot is $\chi_u$ with the
conjugated profile. The other presentations in \eqref{eq-family}
reduce to these by conjugation and finite twists. Test heights
only change sign, so the source profile estimates still apply.

Contour shifts use the full local correction, without dividing
by a possibly vanishing $P_p$. For example the replacement
factor can be written
\[
 G_p=\frac{
  (\overline{\psi(p)}Q^s-Q^{-w})(1-V)(1-W)P_p^*
              -Q^{-w}(1-VW)}{1-D}.
\]
The two Euler regions and the joint error-slot estimates in
Q, Section 16, use only unit-modulus local target values.
Their local algebra is unchanged over $K$. The transferred physical estimates retain the full correction
and its norm-sum estimates, including all its masks.

\subsection{The general-field local factor}

Write
$\vartheta=\bar\Lambda_K\bar G\psi$.  At a good prime of norm $Q$ put
\[
 V=Q^{-6z},\quad W=\chi_u(p)Q^{-w},\quad
 D=\psi(p)\bar\chi_u(p)Q^{-s},\quad
 E=\vartheta(p)^6Q^{4-6s}V.
\]
The exact local sum is \cite[(7.19)]{P},
\[
 P_p=\frac1{1-V}+\frac W{1-W}
 +\frac E{1-E}\left(\frac{1-Q^{-1}}{1-V}+W\right)
 -\frac D{1-E}\left(1+\frac{(Q-1)WV}{1-V}\right).
\]
Thus the full factorization is
\[
 \zeta_K^S(6z)\frac{L_K^S(w,\chi_u)}
                         {L_K^S(s,\psi\bar\chi_u)}H_{u,\psi}(s,w,z).
\]
In particular the sixth-power row variable does not carry $\psi^6$;
both principal poles remain at $w=1$ and $z=1/6$. Define
$H_\psi^S(s)=H_{1,\psi}(s,1,1/6)$, and let $H_{\psi,1,Z}$ denote
the principal-row correction after all compensations and prime-window
sums have been performed.

\subsection{The positive normalizer}

Use the original geometry
\begin{equation}\label{eq-geometry}
 \ell=\frac16,\quad b=\frac18,\quad M=\frac56,
 \quad l_x=\frac{17}{48},\quad l_y=\frac{23}{48},
 \quad h=\frac{13}{16},\quad C(s)=s-\frac{11}{16}.
\end{equation}
Here $X=Z^{l_x}$ and $Y=Z^{l_y}$.
Choose a sufficiently fine fixed even slot count $J_0$, with
$\ell_i=\ell/J_0$, scales $P_i=Z^{\ell_i}$, and nonnegative
tests $W_i$ on disjoint annuli. These analytic tests, lengths,
and the count precede the later target in $\Packet$.
After fixing that target and its arithmetic data, choose the
finite identity-ray filter $T=T_\psi$ needed by the local
presentation, and let $\mathcal P_i$ be the corresponding primes
in the $i$th annulus. Put
\begin{equation}\label{eq-normalizer}
 S_i(Z)=\sum_{p\in\mathcal P_i}W_i(q_p/P_i)q_p^{-5/6},
 \qquad A_T(Z)=(-1)^{J_0} Z^{-\ell/6}\prod_iS_i(Z).
\end{equation}
The fixed-ray prime theorem gives $S_i(Z)>0$ eventually and
$A_T(Z)\asymp_{K,T,W_i}(\log Z)^{-J_0}>0$.
The ray filter changes a fixed positive density and the lower
threshold. It follows that $A_T(Z)^{-1}\ll_{\psi,\eps}Z^\eps$
for every preassigned $\eps>0$. No uniform ray-class density
is required.

At the principal residue $u=1,w=1,z=1/6$, Q's local estimate
is still $B_p=-1+O(q_p^{-7/8})$. Choose, before the target,
\[
 \mu_0=\frac7{16}\min_i\ell_i=\frac7{96J_0}>0.
\]
Then uniformly in $\Im s$ on $\Rea s>7/8$,
\begin{equation}\label{eq-principal-approx}
 H_{\psi,1,Z}(s,1,1/6)
   =H_\psi^S(s)Z^{\ell/6}A_T(Z)
        \{1+O_\psi(Z^{-\mu_0})\}.
\end{equation}
The same fixed set $S$ is used throughout. Enlarging its fixed
cutoff as in Q, (10.1)--(10.2), gives
\begin{equation}\label{eq-contraction}
 H_\psi^S\text{ holomorphic on }\Rea s>7/8,
 \qquad |H_\psi^S(s)-1|\le\tfrac12.
\end{equation}
On the principal residue $u=1,w=1,z=1/6$, the displayed local formula
gives, uniformly for $\sigma=\Re s\ge7/8$,
\[
 H_p-1\ll Q^{-2}+Q^{-1-\sigma}+Q^{3-6\sigma},\qquad
 B_p+1\ll Q^{-\sigma}+Q^{-1}+Q^{3-5\sigma}
             \ll Q^{-7/8}.
\]
The first bound is normally summable.  A fixed enlargement of $S$
makes the logarithmic Euler-product tail as small as desired and
proves \eqref{eq-contraction}.  Once this has been done the quotient
$B_p=G_p/H_p$ is legitimate on this principal locus.
For each fixed positive window, its relative error is at most
$O(P_i^{-7/8})$, by division by the same positive sum $S_i(Z)$.
Thus \eqref{eq-principal-approx} holds with the pretarget
$\mu_0=7/(96J_0)$, even if $T$ was selected later.

We make the base probe normalization explicit, separately from
the prime-slot tests $W_i$. Choose the Bruhat--Schwartz test
$W_{0,\mathrm{add}}$ compactly supported in $K_S^\times$ and the
nonnegative nonzero annular test $W_{1,\mathrm{add}}$, as in
\cite[(7.1)--(7.3)]{P}. Put
\[
 \mathcal M_{\mathrm{add}}(z)
 =\int_{K_S^\times}\widehat W_{0,\mathrm{add}}(y)
                          |y|_S^z\,d^\times y,
 \qquad
 \widehat W_{1,\mathrm{add}}(w)
 =\int_0^\infty W_{1,\mathrm{add}}(t)t^w\,\frac{dt}{t}.
\]
The first test is chosen with
$\widehat W_{0,\mathrm{add}}(0)=0$ and
$\mathcal M_{\mathrm{add}}(1/6)\ne0$.
At one complex place, subtract suitable multiples of two bumps
on different radial subannuli: their additive integrals cancel
while their different radial moments do not. The local zeta
functional equation gives the required nonzero Mellin pairing,
exactly as in P, (7.1). A fixed scalar then makes
$\mathcal M_{\mathrm{add}}(1/6)>0$.
Only $W_{1,\mathrm{add}}$ is required to be nonnegative.

In this normalization P, (8.13), gives the unmarked double-residue
scalar
\[
 c_S=\frac16\bigl(\operatorname*{Res}_{s=1}\zeta_K^S(s)\bigr)^2
         \mathcal M_{\mathrm{add}}(1/6)
         \widehat W_{1,\mathrm{add}}(1)>0.
\]
The Gaussian normalization removes the target-dependent
archimedean scalar before this computation. The physical principal
row is unique: a row character trivial on all exterior primes is
trivial by Chebotarev, so its Kummer class is a sixth power and
its unit ambiguity belongs to $(R^\times)^6$. Any finer unit
sectors used in positive moment estimates are recombined into
this one physical orbit. Compensation commutes with unfolding
and creates no additional principal residue.



With this calibrated positive scalar $c_S$, set
\begin{equation}\label{eq-signal}
 J_\psi(Z)=\frac{I_{\psi,\mathrm{mod}}(Z)}{c_S A_T(Z)},
 \qquad
 f_\psi(Z)=\frac1{2\pi i}\int_{(2)}
     Z^{C(s)}e^{(s-5/6)^2}\frac{H_\psi^S(s)}{L_K^S(s,\psi)}\dd s.
\end{equation}
Neither expression contains an analysis height cutoff.
Lemma~\ref{lem-orbit-gram}, the marked reflected energy in
Appendix~\ref{app-transfer}, and the positive calculation of
Q, Proposition 15.3 give, after the subpower reciprocal normalizer,
\begin{equation}\label{eq-low}
 |J_\psi(Z)|\ll_{\psi,\eps}Z^{3/16+\eps},
 \qquad C(7/8)=3/16.
\end{equation}
No parameter optimization below $7/8$ is needed.


\section{The dynamic endpoint and one family closure}\label{sec-close}

Fix $(K,\omega)$ and suppose $B=B_{K,\omega}>7/8$.
Proposition~\ref{prop-seed} gives $B\le11/12$, so
\eqref{eq-kappa} applies. Every companion denominator in the
detector belongs to $\Packet$, hence its zero real parts are
at most $B$. All common small-row and principal target
contours use $B+e$. This is legitimate before any improved
bound has been proved.

\subsection{The capacity comparison}

Use $\delta=2a-1$, $q=x\delta$, with
$0\le\delta\le\alpha=5/6$ and $0\le x\le1/2$.
At effective row width one, an inverse witness of length $r<1$
has capacity $(1-r)/2$. A plain witness of length $m<1/2$,
used twice in the fourth moment, has the actual capacity
\begin{equation}\label{eq-dynamic-capacity}
 z_P(m)=\frac{1-2m}{6\kappa}
       =\frac{1-2m}{9/2+12\Delta}.
\end{equation}
The baseline $2(1-2m)/9$ is used only for comparison.

For the same witness and spike choices as Q, Proposition 19.2,
the difference in the plain row-count exponent is
\begin{align}\label{eq-capacity-cost}
 2q(1-2m)\left(\frac29-\frac1{9/2+12\Delta}\right)
 &=\frac{24q(1-2m)\Delta}{(9/2)(9/2+12\Delta)}
 \le\frac{\Delta}{4}+O(\eps).
\end{align}
The comparison is used where $m\ge1/3-O(\eps)$ and $2q\le1$.
Zero-capacity endpoints use the independent-scale inverse
amplification or the zero-slot plain moment, rather than a
strict marked estimate at zero width.

Put
\begin{align*}
 D_x&=3-\frac{17x}{9}, & P_x&=\left(2-\frac{8x}{9}\right)(1-x),\\
 J&=(\alpha-\delta)D_x+\delta P_x,
 &R_*&=1-\delta+
       \frac{(\alpha-\delta)\delta P_x}{2J}.
\end{align*}
The crossing of the short and long witnesses in Q, (19.7) and
(20.8), now gives the actual bound
\begin{equation}\label{eq-row-dynamic}
 R\le R_*+\frac{\Delta}{4}+\eps.
\end{equation}
This uses only the inverse and paired plain inputs above.
The witness/row conversion is Q, Sections 8 and 19 applied with
the transferred moments and fixed-type Hecke continuation.
The conductor powers are unchanged; fixed local gamma parameters
affect only height constants.
One has $35/54\le J\le5/2$ on the whole closed parameter rectangle.

The only finite exceptional presentations among the
sixth-power-free physical rows lie in the finite fixed sector
part. Indeed, a prime outside $S$ with valuation from one to
five ramifies the sextic character; a character with fixed
conductor in $S$ has no such prime. Class representatives
and $U_S/U_S^6$ give finitely many remaining sectors. This
assertion concerns physical rows; canonical sixth-power
descendants retain the exceptional treatment in $\mathsf P$.

\subsection{The original numerical certificate}

If $U=Z^d$ is a moderate row norm, the full correction in
$\mathsf T$ gives the relative exponent
\begin{equation}\label{eq-E-general}
 \mathcal E(d;R,q)=
 -\frac1{48}+\frac23\delta+\frac q6
       -h(1-R)
       +(d-h)\left(R+\frac\delta2-\frac{17}{50}\right),
 \qquad h=\frac{13}{16},
\end{equation}
measured against $C(7/8)=3/16$.
The term $q/6$ is already a base-$Z$ exponent.
Changing $R$ by $c$ changes $\mathcal E$ by $dc$.
Write $E_*(d)=\mathcal E(d;R_*,q)$.

The exact certificate Q, (20.9), is unchanged. Set
$y=1/2-x$ and $v=51+41y$. Then
\begin{align}\label{eq-certificate}
 10368vJ(-E_*(h))={}&
 (3+5y)\{(4v\delta-79)^2+49\}\nonumber\\
 &+4y\{4v\delta[(1+3y)(15+32y)\delta+9-13y]
                                      +265+3485y\}.
\end{align}
Every displayed summand is nonnegative for
$0\le y\le1/2$ and $0\le\delta\le5/6$.
Since $v\le17(3+5y)$ and $J\le5/2$,
\begin{equation}\label{eq-high-margin}
 -E_*(h)\ge\frac{49}{440640}>10^{-4}.
\end{equation}
The crucial comparison for the actual dynamic capacity is therefore
\begin{equation}\label{eq-dynamic-saving}
 \mathcal E(h;R,q)-\{C(B)-C(7/8)\}
 \le-\frac{49}{440640}
       +\frac h4\Delta-\Delta+O(\eps)
 =-\frac{49}{440640}-\frac{51}{64}\Delta+O(\eps).
\end{equation}
Thus the loss in capacity is absorbed by the very gap that defines
the contradiction. No fourth moment with $\kappa<3/4$ is used.

For the bounding exponent $R_*+\Delta/4$ the coefficient of
$d-h$ in \eqref{eq-E-general} is at least
$33/50-\delta/2\ge73/300$. Hence the endpoint controls the
selected moderate ranges below $h$. The supply margin is
\[
 \frac{\ell}{h}=\frac8{39}>\frac7{37},\qquad
 \frac8{39}-\frac7{37}=\frac{23}{1443}.
\]
A fixed extension $h<d\le h+\zeta$, $\zeta<1/48$, leaves
available length greater than $1/5>7/37$.
At $\delta=5/6$ the unmarked amplified endpoint gives
$R=1-\delta$; no shrinking strict inverse width is used.

\subsection{The remaining ranges}

Take $d_{\min}=1/100$. The source comparisons retained in
$\mathsf T$ give the following positive margins before small
losses.
\begin{center}
\begin{tabular}{p{.57\linewidth}p{.31\linewidth}}
\toprule
Range & Saving\\
\midrule
Floor bin $\delta=1/50$, before subtracting $\Delta$
& $7/1200$\\
$d_{\min}\le d\le1/2$, no selected slots,
before subtracting $\Delta$ & $49/14400$\\
Small rows $1\le U\le Z^{d_{\min}}$,
apart from the principal row, relative to $C(B)$ & $63/800$\\
Principal $w$ contour & $23/960$\\
Principal $z$ contour & $13/9600$\\
Principal slot approximation & $\mu_0=7/(96J_0)$\\
\bottomrule
\end{tabular}
\end{center}
For the middle range, the unselected count
$R=76/75-2\delta/3$ yields
\[
 \mathcal E(d;R,q)\le
 -\frac{529}{2400}+\frac{25}{96}\delta
 \le-\frac{49}{14400}.
\]
For the small rows the full-correction estimate on
$(\Rea s,\Rea w,\Rea z)=(B+e,1/2,17/50)$ gives
$-79/800+2d_{\min}=-63/800$, relative to $C(B)$.
This includes the nontrivial unit sectors. There is no additional
cost from using $11/12$ instead of an individual target boundary:
the common line is $B+e$ throughout.
Rows above $h+\zeta$ are estimated on
$(2,2,z_\infty)$; the arbitrary-decay assertion in $\mathsf T$
allows a fixed sufficiently large $z_\infty$.

\subsection{Real choices before the target}

The following is Q, Section 20.6's order of choices applied to
the family $\Packet$. Put
\[
 m_{\rm hi}=\frac{51}{64}\Delta,\qquad
 m_{\rm high}=\min\{m_{\rm hi},63/800\}.
\]
Reserve fixed fractions of the low margin $\Delta$, the high
margin $m_{\rm high}$, and both principal contour margins.
Make the choices common for the permitted representatives of
the fixed finite collection of types $0,\omega_\infty,
\bar\omega_\infty$.
Choose the moment losses and capacity decrements first.
The uniform real-parameter assertion of $\mathsf P$ then
gives a slot mesh on $3/4\le\kappa\le5/6$.
Choose an even $J_0$ sufficiently large that
$2\ell/J_0$ is below that mesh, the allocated rounding loss,
and $1/185$.

Take $\ell_i=\ell/J_0$ and choose nonnegative nonzero tests on
disjoint intervals
\[
 I_i=\left(1+\frac{2i-1}{2J_0+1},
           1+\frac{2i}{2J_0+1}\right),\qquad 1\le i\le J_0.
\]
These analytic supports are disjoint before any ray or row mask
and are chosen before the finite-order target twist. The fixed
arithmetic ray filter may be selected after that twist; it changes
only constants and thresholds. The pretarget error exponent is
$\mu_0=7/(96J_0)$ in \eqref{eq-principal-approx}.
Now choose the amplitude and bin widths, remaining real
losses, $\zeta$, and $z_\infty$, in this order, reserving
the fractions used in Q, Section 20.6.
The real-parameter uniformity for these choices follows from the
transfer and finite-depth argument in Appendix~\ref{app-transfer}.

For definiteness, after these allocations one can retain
\begin{equation}\label{eq-sigma}
 \sigma=\frac18\min\left\{
      \frac{51}{64}\Delta,\frac{63}{800},
      \frac{23}{960},\frac{13}{9600},\frac7{96J_0}
                       \right\}>0.
\end{equation}
This number depends on the fixed family and $\Delta$, but
not on its later target $\psi$. Notice that
$\sigma<\Delta/2$.

\subsection{Height choices and continuation}

Now fix a target $\psi\in\Packet$, removing its global norm
twist. Choose between it and its conjugate to satisfy the
nonzero angular condition in Section \ref{subsec-general-arch}.
Conjugation preserves the family and the real parts of zeros,
so proving the estimates for this representative suffices.
Choose its arithmetic data, final $S$, finite sectors, and
identity-ray filter $T$. Then choose all retained internal
seminorm and Sobolev orders.
The transferred estimates give finite numbers
$A_\psi,B_\psi$ and a positive permitted height ceiling such
that, for the subsequently chosen external order $N$,
\[
 |J_\psi-f_\psi|
 \ll_{\psi,N}
 Z^{C(B)-2\sigma}(1+T_1)^{A_\psi}
       +Z^{B_\psi}T_1^{-N}.
\]
Degree, infinity type, and fixed local profiles can enlarge
$A_\psi$, $B_\psi$, and the threshold. Choose a positive
$\tau_\psi$ small enough for the height ceiling and the
reserved saving, put $T_1=Z^{\tau_\psi}$, and then choose
$N$ large enough. With the strict reserves above, the
conclusion of Q, Lemma 11.1 is
\begin{equation}\label{eq-high-final}
 |J_\psi(Z)-f_\psi(Z)|\ll_\psi Z^{C(B)-\sigma}.
\end{equation}
The analysis height belongs only to this estimate:
neither $J_\psi$ nor $f_\psi$ depends on $T_1$.
The low bound with loss $\Delta/2$ and
$\sigma<\Delta/2$ now gives
\begin{equation}\label{eq-f-bound}
 |f_\psi(Z)|\ll_\psi Z^{C(B)-\sigma}.
\end{equation}

The continuation is elementary once this common exponent is
available. Moving the defining line of $f_\psi$ arbitrarily far
right gives rapid decay as $Z\downarrow0$. Hence
\[
 \mathcal F_\psi(s)=\int_0^\infty
       f_\psi(Z)Z^{-C(s)}\,\frac{dZ}{Z}
\]
is holomorphic on $\Rea s>B-\sigma$.
Mellin inversion, first on $\Rea s=2$, identifies it with
\[
 e^{(s-5/6)^2}\frac{H_\psi^S(s)}{L_K^S(s,\psi)}.
\]
Both numerator factors are nonzero on $\Rea s>7/8$, and
$B-\sigma>7/8$. Thus every $L_K(s,\psi)$ in $\Packet$
has no zero in $\Rea s>B-\sigma$. At a norm-character
pole its reciprocal merely has a zero, which causes no
problem.

By the definition of $B$, if $B>7/8$ there exists a target
in $\Packet$ with a zero of real part greater than
$B-\sigma/2$. The preceding conclusion is a contradiction.
Consequently $B\le7/8$, and in particular
$L_K(s,\omega)$ is zero-free in $\Rea s>7/8$.
Apply this with $\omega=\psi\circ N_{K/F}$ and use
Lemma \ref{lem-descent}. This proves Theorem \ref{thm-main}.

\section{Original sources and proof dependencies}\label{sec-scope}

The main theorem concerns arbitrary number fields and arbitrary
fixed unitary infinity types. Its analytic source dependence is
Source Input~\ref{hyp-source}. The general-field conclusions
$\mathsf T$, $\mathsf I$, and $\mathsf P$ are the proved transfer
of Proposition~\ref{hyp-main}, not further hypotheses.

\begin{center}
\begin{tabular}{>{\raggedright\arraybackslash}p{.45\linewidth}
                >{\raggedright\arraybackslash}p{.46\linewidth}}
\toprule
Source or argument & Role in this paper\\
\midrule
P, Sections 3, 5--7 and Appendix A
& Original theta construction, arithmetic calibration, and local
coefficient identities; its general-field and fixed-type
realizations are verified in Section~\ref{sec-inputs}.\\
Q, Sections 5--16
& Original physical and marked-reflection constructions;
norm-form Poisson, additive Gram, and full local compensation
are verified here.\\
Q, Sections 17--18
& Original long canonical proofs. Appendix~\ref{app-transfer}
supplies the general-field replacements, keeping their integer
width and capacity calculations.\\
Q, Sections 19--20
& Witness comparison and the exact endpoint certificate;
the dynamic capacity loss and fixed-type family closure are
calculated in Section~\ref{sec-close}.\\
Abelian descent and the fixed-type starting bound
& Proved here, including norm-character poles and exact
sixth-power amplification masks.\\
\bottomrule
\end{tabular}
\end{center}

Q's published finite-order $7/8$ theorem and P's cyclotomic
$1-10^{-6}$ theorem are not used as a shortcut to an arbitrary-field
conclusion. In particular $K=F(\zeta_{12})$ need not be cyclotomic
over $\mathbb Q$. The transfer is the reason the original
constructions apply to this $K$.

The higher-order large sieve of \cite[Theorem 1.3]{BGL} is compatible
with the exterior ideal symbols, but it does not alone prove the
canonical inverse moment. At column length $N=M^{9/10}$ its term
$(MN)^{2/3}$ is $M^{19/15}$, exceeding the desired $M^{1+\eps}$.
The two source Poisson reductions and the marked terminal are
therefore essential. Likewise, the actual identity
$\chi_n(a^6)=\one_{(n,a)=1}$ is retained throughout; replacing it
by a primitive character and discarding its zeros would invalidate
the amplification.

This paper reproduces the transfer arguments and the simplified
closure. It does not independently re-prove every line of the
long original theta and canonical constructions. That is the
remaining source-relative qualification in Theorem~\ref{thm-main}.

\appendix
\section{Transfer of the canonical estimates}\label{app-transfer}

\begin{proposition}[Fixed-field moment transfer]
\label{prop-recovered-moment-transfer}
Retain the original canonical arguments of Q, Sections 14, 17 and 18,
and the normalized theta realization of P, with its general-field and
fixed-infinity-type realization established above. Then the directed
marked inverse estimate \eqref{eq-MI} and the paired fourth moment
\eqref{eq-fourth} hold over the fixed field $K$ with their displayed
width, coefficient, zero-extension and common-mask restrictions.
For the fourth moment one may use $3/4\le\kappa\le1$; when
$\kappa<1$, its prime premise concerns only finite-order characters
and is the boundary $(1+\kappa)/2$.
The slot mesh is independent of the subsequently fixed slot count.
For a fixed infinity-type orbit, the real losses and that mesh can
also be chosen before the later finite-order target twist and its
fixed conductor. Thus the two general-field moment estimates are
consequences of the cited canonical proofs and the replacements
verified here, rather than additional moment hypotheses.
\end{proposition}

\begin{proof}\leavevmode\par
\paragraph{1. The common orbital Poisson formula.}
Use Lemma~\ref{lem-orbit-gram}, including its general-periodic
version and separate zero orbits. The primitive character formula
is used in the two inverse transforms and their initialization;
the full finite Fourier formula is used in the second plain
transform. The mask is expanded first. Its common norm kernel,
fixed logarithmic growth, and finite-seminorm bounds are exactly
the profile data required by the source reductions.

\paragraph{2. Local tables, masks, and fixed unitary coefficients.}
At a good prime of norm $P$ the full normalized finite sums obey
\[
 |G(\mathfrak p^a,h)|=P^{(a-1)/2}
       \one_{v_{\mathfrak p}(h)=a-1}\quad(6\nmid a),
 \qquad
 G(\mathfrak p^a,h)=P^{a/2}\one_{\mathfrak p^a\mid h}
       -P^{a/2-1}\one_{\mathfrak p^{a-1}\mid h}\quad(6\mid a).
\]
These follow by separating the higher residue digits and applying
finite-field orthogonality. CRT gives the full two-column
correlations, including collision zeros, exactly as in Q,
Section 13.3. Thus a displayed sixth-divisible exponent remains
a puncture. All ideal common-support extractions are exact before
any coefficient is bounded or a coprime support is enlarged.
Once factor ideals and the finitely many unit sectors are fixed,
the old-label fibre is fixed modulo $\Gamma$; this replaces the
source's uniqueness of primary generators, at a fixed finite
cost. The fixed background stays on columns, not in an exterior
row modulus. Lemma \ref{lem-directed} gives its exact propagation:
the first and second transforms replace $\nu$ only by
$\nu\vartheta$, $\vartheta\in\Theta$, using
$\nu(ACn)=\nu(A)\nu(C)\nu(n)$ and $|\nu(C)|^2=1$.
No finite Fourier expansion of a genuine $\nu$ is made.
Character-equivariant theta periods retain the fixed unit phases;
their conductor phases are multiplicative row, tuple, or frozen
scalars. They are not dropped by replacing the period by an
unweighted unit average. Cube completion leaves $\nu(b)^3$ on
the extracted cube variable, as already checked in that lemma.

\paragraph{3. The full marked terminal.}
The identity
$\one_{p\mid A}=1-\chi_p(A)^0$ applied to the compatible
theta branches gives inactive and active mark sizes $P^{-1}$
and $P^{-1/2}$. A residual row digit $j=1$ and a mark digit
$j=0$ have trivial pair phase. The remaining row character is
quadratic and the active marked squarefree column is cubic.
Freeze every other digit, including $j=4$, and separate fixed
local sectors of the row product and marked product individually.
The annular theta kernel established above, the fixed-field quadratic sieve verified in
\cite[Lemma 6.4]{P}, and the cubic sieve of \cite[Theorem 1.3]{BGL}
therefore give Q, Lemma 14.3's full reflected energy. In its
notation this exponent is
\[
 E_{\rm ref}=\frac O2+\max(H,v+\ell_b)-S_0-B_0+z_a-u-\ell_b
       -\frac12(T_d-v-3\ell_b)_+,
 \qquad u=\min\{v,z_a,(v+z_a)/3\}.
\]
Here $O$ is the powerful-row length, $H$ the actual residual
squarefree-row length, $v,\ell_b$ the dual squarefree and cube
lengths, and $S_0,B_0,T_d$ are the local extraction quantities of
that lemma. Its completed quadratic reduction, Q, Lemma 5.5,
is an ideal factorization with all collision masks retained.
No special low-side condition $2A_0\le O$ is used here.
For the squarefree fourth-power label of length $V$, puncture
length $Q$, row length $M$, and completed scale $N_*$, the
primewise conductor comparison is
\[
 H\le M-O+o(1),\qquad
 T_d-S_0-B_0\le2M-O+Q+V-N_*+2z_a+o(1).
\]
These are Q, (17.11), with ideal norms. If $F_0=N+V$ and
\[
 F_0-M-Q-z_0\ge c,\qquad 4F_0-3M-6z_0\ge c,
\]
the short completion $V<d_0$, $N_*=N-3b_h$,
$b_h<d_0+o(1)$ gives the three terminal bounds
\[
 F_0-c,\qquad F_0-2c+5d_0,\qquad
 F_0-M/4-3c/4+5d_0/2,
\]
apart from the assigned small losses. They are strictly below
$F_0$ when $d_0\le c/200$, precisely Q, (17.12)--(17.14).
The label mass is counted once and cancels its original
$Z^{-V}$ normalization. Neither a label prime nor a puncture
prime has been incorporated into the fixed data.

\paragraph{4. The two inverse transforms and initialization.}
Use the orbital identity in paragraph 1 at Q, (17.22) and
(17.36), retaining the genuine coprime expression before
truncating its ratio tails. Its prefactor, conductor norm,
puncture divisor, finite transform and scalar kernel argument
are exactly the quantities in Q, (17.23)--(17.60).
Paragraph 2 gives the same finite column identities and
whole-product priority assignments. In particular the canonical
coefficient remains
$\overline\Lambda\gamma_2\nu\rho$, since the calibrated
squarefree identities are
$\gamma_2^3=\mu_K\Lambda$ and
$\overline\Lambda\gamma_2\gamma_1=\mu_KG$.
The two invariants in paragraph 3 and the strict child-row
decrease in Q, (17.17), are unchanged. The common smooth kernel
and all outside row masks remain common through the weighted
Cauchy inequalities of Q, Section 17.3; no row-dependent column
coefficient is introduced. The one-transform initialization has
the same local identities. Its two margins, before enclosure
losses, are
\[
 m-r-2z+2G,\qquad 3m-2r-8z+10G+2P_1,
\]
as in Q, (17.85)--(17.86). Here $G,P_1\ge0$ are the source's
extracted norm lengths. The hypotheses of \eqref{eq-MI}
supply the required positive margins. This imports the finite
canonical induction and proves the directed inverse estimate,
with its exact mask and coefficient restrictions. Proposition~\ref{prop-seed} gives the independent-scale
form with the exact sixth-power masks.

\paragraph{5. The ordinary fourth-moment reflection and coefficients.}
The compatibility checks in Section \ref{sec-fourth} prove the
fixed-infinity-type reflection and the paired propagation, so
we use them rather than treating an infinite-order twist as a
finite Fourier coefficient. The finite conductor factor remains
$Q^{1/2-s}$ and the dual norm scale $Q/X$. The product of complex
gamma quotients may have repeated poles at real part zero;
its small annular coefficients are
$O(2^{-j/2}(1+j)^A)$, which are summable. The degree changes
orders and constants, not the conductor exponent.
The common-mask identities Q, (18.4)--(18.8), are applied to
the displayed presentation before primitive reduction.
Apply the general-periodic version of paragraph 1 in the two
plain transforms. CRT and complete factor allocation preserve
Q, Lemma 18.2's one whole-product character on each separated
side, the original individual slot coefficients, one common
mask and norm power on both rectangles, and their equal product
of plain scales. The two sides differ by a character in
$\Theta$. Thus they have the same finite exceptional status,
and the actual prime characters remain finite order. Their
prime estimate is Q, (18.12), with the stated finite-order
zero-free premise. All strict width and positive-part ledgers
in Q, Sections 18.4--18.8, consequently keep their coefficients.

\paragraph{6. Exceptional rows and the prime pool.}
The unmasked Mellin argument in Section \ref{sec-fourth},
followed only afterwards by common-mask inclusion--exclusion,
replaces the exceptional lattice estimate. If the four plain
scales are at least $Z^a$, $a\ge0$, and have common product
$T_0$, its equal-product cancellation gives
$O(Z^\eps T_0^{1/2}Z^{-a}(1+|t|)^A)$ after central
normalization, exactly the saving in Q, (18.47); a genuine
background has zero residue and gives a stronger bound.
The canonical exceptional condition prescribes each exterior
row valuation modulo six. Hence its ideal is
$\mathfrak h_0\mathfrak a^6$ in finitely many unit sectors.
At a newly nonunit extracted prime of common multiplicity one,
the fixed factor has valuation two, so the row valuation is
four modulo six. If these primes have length $v_1$, then
$\N\mathfrak h_0\ge Z^{4v_1}$, giving the count required by
Q, (18.41). The old moving support is disjoint from this
support because its original column zeros are still present.
The comparison of $h$ with $hp^6$ uses a fixed prime pool of
length $\ell_{\rm pool}$; the prime ideal theorem supplies
$Z^{\ell_{\rm pool}+o(1)}$ allowed primes. An output ideal
has only a bounded number of such prime divisors. The same
single pool-density gain therefore survives on unit orbits.
The local error valuations $1,6,7$ have squared coefficients
$P^{-1},1,P^{-1}$ by paragraph 2. A mesh below
$\ell_{\rm pool}/2$ separates the pool from all live slots.
These are the field-sensitive ingredients in the plain
induction, and prove \eqref{eq-fourth} by the cited induction.

\paragraph{7. Uniform real choices and finite propagation.}
The source chooses its width floor, strict decrements and
finite depth from bounded real length ranges and requested
losses. None of the replacements above changes a norm exponent
or an affine ledger. A fixed $K,S,\psi$ can add finite sectors,
kernel logarithmic powers, seminorm orders and finite constants.
For every preassigned $\eta>0$, each such logarithm is bounded
by $C_{K,S,\psi,\eta}Z^\eta$; the same $\eta$ works whatever
its fixed logarithmic degree. Allocate this aggregate loss
through the source's finite depth before choosing the internal
orders. Support-enclosure constants affect only the eventual
threshold through $O_{K,S,\psi}(1/\log Z)$.
The width step fixes the prime-pool length, and then the mesh,
before the number of slots is fixed. Divisor allocations for
that fixed number of slots cost their one small-power loss;
their constants, not the mesh, may depend on the slot count.
It follows that real choices can precede the later finite-order
twist and enlargement of $S$. Their constants and height orders
need not do so. The common Fourier measures in paragraphs 1
and 4 allow the finite-seminorm propagation of Q, Section 17.3
and the loss order of Section 18.8 unchanged. Increasing a final
external tail order differentiates only external tests, not an
already chosen internal moment profile, so its height degree
does not increase. Finer unit sectors used in positive moment estimates are recombined
into the original $(R^\times)^6$ physical orbits; they add no
principal residue. The analytic annuli, slot lengths, and number of slots precede
the target conductor. The fixed arithmetic ray filter can be
chosen afterwards: it changes the positive normalizer by a
fixed density, while the relative local slot error and its
preassigned exponent $\mu_0$ are unchanged. This proves the asserted
quantitative and mesh uniformity and completes the transfer.
\end{proof}


\begin{thebibliography}{9}
\bibitem{Q} OpenAI,
\emph{The Quasi-Riemann Hypothesis: A Zero-Free Half-Plane
$\operatorname{Re}(s)>7/8$}, 30 September 2026.
\href{https://github.com/openai/math/blob/main/preprints/The-Quasi-Riemann-Hypothesis-September-30-2026/paper.pdf}{Original paper}.

\bibitem{P} OpenAI,
\emph{Primitive roots for every admissible integer base},
4 October 2026.
\href{https://github.com/openai/math/blob/main/preprints/Primitive-roots-for-every-admissible-integer-base-October-4-2026/primitive-roots-all-integer-bases.pdf}{Original paper}.

\bibitem{Tate} J. T. Tate,
\emph{Fourier analysis in number fields and Hecke's zeta-functions},
in J. W. S. Cassels and A. Fr\"ohlich (eds.),
\emph{Algebraic Number Theory}, Academic Press, 1967,
pp. 305--347.

\bibitem{BGL} V. Blomer, L. Goldmakher, and B. Louvel,
\emph{L-functions with $n$-th-order twists},
International Mathematics Research Notices 2014, no. 7,
1925--1955.
\href{https://doi.org/10.1093/imrn/rns257}{doi:10.1093/imrn/rns257};
\href{https://arxiv.org/abs/1112.1650}{arXiv:1112.1650}.
\end{thebibliography}
\end{document}
